\chapter{Fundamentação teórica}

Ao longo deste capitulo será descrito topicos importantes para o entendimento deste trabalho.
Os conceitos apresentados serão: (i) características de qualidade em código, (ii) cobertura de testes, (iii) complexidade ciclomática,
(iv) complexidade cognitiva, (v) análise estática e (vi) a ferramenta SonarQube, que foi utilizada no processo de avaliação.

\section{Características de qualidade em código}

De acordo com a ISO/IEC 9126, a qualidade de software pode ser avaliada com base em características como
funcionalidade, confiabilidade, usabilidade, eficiência, manutenibilidade e portabilidade. Todos esses atributos tornam um codigo
mais agradáveis para desenvolvedores e usuários, além de menos propenso a falhas.

Dentro desses podemos entrar ainda mais a fundo em cada um 

Apresentar os atributos que tornam um código mais legível, manutenível,
consistente e menos propenso a falhas. Nesta seção, também devem ser
abordadas métricas que ajudam a observar a qualidade do código de forma
objetiva.

\subsection{Cobertura de testes}

A cobertura de testes indica a proporção do código exercitada pelos testes
automatizados. Ela não garante ausência de defeitos, mas ajuda a identificar
partes pouco verificadas e a orientar a ampliação da suíte de testes.

\subsection{Complexidade ciclomática}

McCABE (1976) desenvolveu uma métrica para medir a complexidade de módulos,
ou seja, unidades de código fonte de software. Tal métrica foi nomeada complexidade
ciclomática e é amplamente utilizada, de acordo com MICHURA e CAPRETZ (2005). Uma
alta complexidade ciclomática significa que determinada unidade de software é difícil de se
entender e dar manutenção. A métrica é baseada no conceito de caminhos que um código
fonte pode seguir. Quanto mais caminhos existirem, maior é a dificuldade de um
desenvolvedor conseguir abstrair a execução em seu pensamento.

A complexidade ciclomática mede o número de caminhos independentes em um
módulo. Valores elevados tendem a indicar trechos com maior dificuldade de
compreensão, manutenção e teste.

\subsection{Complexidade cognitiva}

A complexidade cognitiva procura refletir o esforço mental necessário para
entender um trecho de código. Diferente de outras métricas estruturais, ela
valoriza a legibilidade e a simplicidade da lógica implementada.

\section{Análise estática}

Explicar o papel da análise estática na identificação de problemas sem
executar o software, destacando ganhos de confiabilidade e produtividade.

\section{SonarQube}

Descrever a ferramenta SonarQube, suas métricas, regras e como ela pode ser
usada para apoiar a avaliação da qualidade do código, especialmente na
verificação de cobertura de testes, complexidade ciclomática e complexidade
cognitiva.
